\section{The Fathers after Dionysius}\label{sec:fathers-after}

From Gregory the Great to the Carolingian schools the Latin Church received
its teaching on the angels chiefly from one pope. Gregory preached the nine
orders to the people of Rome, named the office proper to each, explained why
Michael, Gabriel, and Raphael bear the names they bear, and set the fall of
the first angel and the confirmation of the rest at the centre of his long
exposition of Job. Isidore of Seville gathered Gregory's teaching and
Augustine's into the definitions and sentences that the Latin schools copied
for centuries; the Venerable Bede carried both into his commentaries on Luke
and the Apocalypse and into the history of the English Church, where angels
meet the dying and carry the souls of the just; and Rabanus Maurus handed
Isidore's chapter and Gregory's homily together to the schools of the
Carolingian age. From these witnesses the West learned that ``angel'' is the
name of an office and not of a nature, that the heavenly city is built of
angels and men, that the elect are gathered into the choirs according to the
likeness of their lives, that some spirits stand before God while others are
sent, that the prince of the fallen was the first of the works of God and
fell by pride, and that the angels who stood were made firm by his fall.
Aquinas draws on them in his questions on the bodies of angels, their
knowledge, their fall, their speech, their orders, their government of the
world, their missions, and their guardianship (for the Greek line after
Dionysius, with John Damascene, see \S\ref{sec:byzantine}).

\sectionguard
\subsection{Gregory the Great}\label{sec:fa-gregory}

Gregory, a monk before he was bishop of Rome and a preacher while he
governed the Roman Church, left three works in which the angels stand at the
centre. The forty \emph{Homilies on the Gospels}, preached to the Roman
people in the basilicas of the city, include the sermon on the lost sheep and
the lost drachma in which he laid out the nine orders of angels, their names,
and their ministries. The thirty-five books of the \emph{Morals on Job},
addressed to Leander, bishop of Seville, expound the angels who stand before
God and are sent, the prince of the nations, the number of the heavenly
host, and the pride and ruin of Behemoth and Leviathan, whom Gregory reads as
the first of the angels. The fourth book of the \emph{Dialogues} gathers
the testimonies of his own Roman world to the life of the soul after death,
and there the angels come to the dying and receive the souls of the
just.\footnote{The Latin of the \emph{Homilies on the Gospels} follows the
text transcribed on Latin Wikisource, checked against H.~Hurter's edition
(Innsbruck, 1892), cited by homily and section. The \emph{Morals} are quoted in the Oxford
translation of the Library of the Fathers (1844--1850), cited by book,
chapter, and section, or by book and section (\S) where the translation marks no chapter. The \emph{Dialogues} are quoted in the English of
1608 as re-edited by E.\,G. Gardner (London, 1911), whose chapter numbers are
used.} Aquinas cites Gregory more often than any Latin Father but Augustine
in his questions on the angels, and his account of the orders
(\ST{I q.~108 aa.~5--6}) is a comparison of Gregory and Dionysius.

\subsubsection{The lost sheep and the ten drachmas}

Gregory preached the thirty-fourth homily in the basilica of Saints John and
Paul on the third Sunday after Pentecost, on the parables of the lost sheep
and the lost drachma (Luke 15:1--10). He reads the shepherd as the Creator
of both natures: \latin{Quia enim centenarius perfectus est numerus, ipse
centum oves habuit cum angelorum substantiam et hominum creavit.} A hundred
is the perfect number, and the Lord had a hundred sheep when he created the
substance of angels and of men. One sheep perished when man, by sinning,
left the pastures of life; the ninety-nine left in the desert are \latin{illos
summos angelorum choros}, the highest choirs of the angels, left in heaven.
Heaven is called a desert because it was deserted: man abandoned it when he
sinned. The number of the rational creature made to see God was lessened
when man perished, \latin{et ut perfecta summa ovium integraretur in coelo,
homo perditus quaerebatur in terra}---and that the perfect sum of the sheep
might be made whole in heaven, lost man was sought on earth (\emph{Hom.\ in
Ev.}\ 34.3). The shepherd lays the sheep on his shoulders because in
taking human nature he bore our sins, and he goes home because, man being
restored, he returned to the heavenly kingdom. There he finds his friends and
neighbours, the choirs of angels: friends, \latin{quia voluntatem ejus
continue in sua stabilitate custodiunt}, because they keep his will
continually in their steadfastness; neighbours, \latin{quia claritate
visionis illius sua assiduitate perfruuntur}, because they enjoy the
brightness of his vision without ceasing. He does not say \latin{Congratulamini inventae ovi}, rejoice with the
sheep that is found, but \latin{Mihi}, with me, \latin{quia videlicet ejus
gaudium est vita nostra}: our life is his joy, and when we are brought back
to heaven we fill up the solemnity of his gladness (34.3).

The joy ``before the angels of God upon one sinner doing penance'' (Luke
15:10) is the joy of the captain who loves the soldier who fled and returned
to press the enemy hard more than the one who never turned his back and never
did a brave deed (34.4); yet the just who weep over sins of thought as others
weep over sins of deed give heaven a joy that no penance of sinners
surpasses (34.5).

The woman who lost one of her ten drachmas is the Wisdom of God. The drachma
bears an image, and man, made to the image of God, lost it when he withdrew
by sin from the likeness of his Maker. The woman lights a lamp because the
Wisdom of God appeared in humanity: the lamp is light in an earthen vessel,
the Godhead in the flesh. When she has found the drachma she calls her
friends and neighbours, \latin{illae potestates coelestes~\dots\ quae tanto
supernae sapientiae juxta sunt, quanto ei per gratiam continuae visionis
appropinquant}---the heavenly powers, near to the Wisdom above in the measure
that they approach it by the grace of continual vision. And the ten drachmas
declare the number of the rational creation:

\begin{quote}
\latin{Angelorum quippe et hominum naturam ad cognoscendum se Dominus
condidit, quam dum consistere ad aeternitatem voluit, eam procul dubio ad
suam similitudinem creavit. Decem vero drachmas habuit mulier, quia novem
sunt ordines angelorum. Sed ut compleretur electorum numerus, homo decimus
est creatus, qui a conditore suo nec post culpam periit.} (\emph{Hom.\ in
Ev.}\ 34.6)
\end{quote}

\noindent The Lord made the nature of angels and of men to know him, and
since he willed it to endure for eternity he created it to his likeness. The
woman had ten drachmas because there are nine orders of angels, and man was
created the tenth so that the number of the elect might be complete; nor did
man perish to his Creator even after his fault, for eternal Wisdom, shining
with miracles in the flesh, restored him by the light of the lamp.

\subsubsection{The nine orders gathered from Scripture}

Gregory then shows from Scripture that the orders are nine:

\begin{quote}
\latin{Novem vero angelorum ordines diximus, quia videlicet esse, testante
sacro eloquio, scimus angelos, archangelos, virtutes, potestates,
principatus, dominationes, thronos, cherubim, atque seraphim.}
(\emph{Hom.\ in Ev.}\ 34.7)
\end{quote}

\noindent We have said there are nine orders because, on the witness of Holy
Writ, we know that there are angels, archangels, virtues, powers,
principalities, dominations, thrones, cherubim, and seraphim. Angels and
archangels appear on almost every page of Scripture, cherubim and seraphim in
the prophets; Paul names four orders to the Ephesians, ``above all
principality and power and virtue and dominion'' (Eph 1:21), and to the
Colossians adds the thrones, ``whether thrones, or dominations, or
principalities, or powers'' (Col 1:16). Five orders named by the Apostle and
four by the rest of Scripture make nine: \latin{procul dubio novem esse
angelorum ordines inveniuntur} (34.7).

The prophet confirms the number in his lament over the first angel. To him it
is said, ``Thou wast the seal of resemblance, full of wisdom, and perfect in
beauty'' (Ezek 28:12), and Gregory weighs the phrase:

\begin{quote}
\latin{Ubi notandum quod non ad similitudinem Dei factus, sed signaculum
similitudinis dicitur, ut quo in eo subtilior est natura, eo in illo imago
Dei similius insinuetur expressa.} (\emph{Hom.\ in Ev.}\ 34.7)
\end{quote}

\noindent He is not said to be made to the likeness of God, but to be the
seal of the likeness, so that the subtler his nature, the more closely the
image of God is shown impressed upon him. The prophet adds at once, ``every
precious stone was thy covering: the sardius, the topaz, and the jasper, the
chrysolite, and the onyx, and the beryl, the sapphire, and the carbuncle, and
the emerald'' (Ezek 28:13). \latin{Ecce novem dixit nomina lapidum, quia
profecto novem sunt ordines angelorum.} He named nine stones because there
are nine orders of angels; and the first angel was adorned and covered with
them, \latin{quia dum cunctis agminibus angelorum praelatus est, ex eorum
comparatione clarior fuit}, because, being set over all the hosts of angels,
he was the brighter by comparison with them (34.7). Aquinas makes this
sentence the authority of his \emph{sed contra} when he asks whether the
highest angel who sinned was the highest of all, and he judges Gregory's
answer ``the more probable view'' against Damascene's opinion that the chief
of the fallen was set over the earthly order (\ST{I q.~63 a.~7}; \S\ref{sec:q63}).

Gregory's list ascends: angels, archangels, virtues, powers, principalities,
dominations, thrones, cherubim, seraphim. Dionysius, whom Gregory knew and
names in the same homily, sets the virtues between the dominations and the
powers, and the principalities between the powers and the archangels.
Aquinas determines that the two orderings agree ``as regards all except the
`Principalities' and `Virtues,'\,'' that each ``may claim authority from the
words of the Apostle''---Dionysius from Ephesians, which ranks the virtues
between the powers and the dominations, Gregory from Colossians, which sets
the principalities between the dominations and the powers---and that ``little
or no difference exists in reality'' between them, since Gregory's virtues,
who work miracles, answer to the Dionysian principalities, and the
Dionysian virtues, who give strength to the lower spirits, to Gregory's
principalities (\ST{I q.~108 a.~6} and ad~4; \S\ref{sec:q108};
Appendix~\ref{app:orderings}).

\subsubsection{The name of an office}

Having counted the choirs, Gregory turns to their ministries:

\begin{quote}
\latin{Graeca etenim lingua angeli nuntii, archangeli vero summi nuntii,
vocantur. Sciendum quoque quod angelorum vocabulum, nomen est officii, non
naturae. Nam sancti illi coelestis patriae spiritus semper quidem sunt
spiritus, sed semper vocari angeli nequaquam possunt, quia tunc solum sunt
angeli, cum per eos aliqua nuntiantur.} (\emph{Hom.\ in Ev.}\ 34.8)
\end{quote}

\noindent The holy spirits of the heavenly country are always spirits, but
they are angels, messengers, only when something is announced through them;
so the Psalmist says, ``Who makest thy angels spirits'' (Ps 103[104]:4), that
is, those whom he always has as spirits he makes angels when he wills. Those
who announce the least things are angels, those who announce the highest are
archangels, and so the archangel Gabriel was sent to the Virgin Mary (Luke
1:26): \latin{Ad hoc quippe ministerium summum angelum venire dignum fuerat,
qui summum omnium nuntiabat}---it was fitting that the highest angel should
come to this ministry, since he announced the highest of all things (34.8).
The Catechism teaches the same distinction in Augustine's words: ``angel'' is
the name of their office, ``spirit'' the name of their nature (CCC~329).

The same rule governs the proper names. The angels do not receive names in
the holy city as though their persons could not be known there without them,
for that city is made perfect by full knowledge from the vision of almighty
God; \latin{sed cum ad nos aliquid ministraturi veniunt, apud nos etiam
nomina a ministeriis trahunt}---but when they come to us to minister, they
draw their names among us from their ministries (34.8).

\begin{quote}
\latin{Michael namque, quis ut Deus; Gabriel autem, fortitudo Dei; Raphael
vero dicitur medicina Dei.} (\emph{Hom.\ in Ev.}\ 34.9)
\end{quote}

\noindent Whenever a work of wonderful power is done, Michael, Who is like
God, is said to be sent, so that the deed and the name together teach that
no one can do what God can do. The ancient enemy who said, ``I will ascend
into heaven,~\dots\ I will be like the most High'' (Isa 14:13--14), will at
the end of the world fight with the archangel, as John saw: ``Michael and his
angels fought with the dragon'' (Apoc 12:7).

\begin{quote}
\latin{ut qui se ad Dei similitudinem superbus extulerat, per Michaelem
peremptus discat, quia ad Dei similitudinem per superbiam nullus exsurgat.}
(\emph{Hom.\ in Ev.}\ 34.9)
\end{quote}

\noindent He who proudly raised himself to the likeness of God will learn,
slain by Michael, that no one rises to God's likeness by pride. Gabriel, the
strength of God, is sent to Mary because \latin{illum quippe nuntiare
veniebat, qui ad debellandas aereas potestates humilis apparere dignatus
est}: he came to announce him who deigned to appear humble in order to
overthrow the powers of the air, ``the Lord mighty in battle'' and ``the Lord
of hosts'' (Ps 23[24]:8, 10). Raphael, the medicine of God, touched the eyes
of Tobias and wiped away the darkness of his blindness (Tob 11; 34.9). The Congregation for Divine
Worship, in its \emph{Directory on Popular Piety and the Liturgy} (2001),
discourages assigning names to the holy angels other than these three, which
Holy Scripture gives (no.~217; \S\ref{sec:devotion}).

Gregory then runs through the names of the other orders, each taken from its
office (34.10). The \emph{virtues} are those spirits \latin{per quos signa
et miracula frequentius fiunt}, through whom signs and miracles are most
often done. The \emph{powers} have received this more mightily than the
rest of their order, \latin{ut eorum ditioni virtutes adversae subjectae
sint, quorum potestate refrenantur, ne corda hominum tantum tentare
praevaleant quantum volunt}: the adverse powers are subject to their rule and
are curbed by their power, lest they tempt the hearts of men as much as they
wish. The \emph{principalities} preside over the good spirits themselves and,
disposing what is to be done, lead the others in fulfilling the divine
ministries. The \emph{dominations} surpass the principalities and powers by
a lofty difference, and Gregory distinguishes the two words with care:
\latin{Nam principari est inter reliquos priorem existere, dominari vero est
etiam subjectos quosque possidere}---to preside is to be first among the
rest, to dominate is to possess those who are subject. The \emph{thrones}
are the hosts over whom almighty God always presides to execute judgment;
since ``throne'' in Latin is \latin{sedes}, seat, \latin{throni Dei dicti
sunt hi qui tanta divinitatis gratia replentur, ut in eis Dominus sedeat, et
per eos sua judicia decernat}: those are called God's thrones who are filled
with such grace of the Godhead that the Lord sits in them and decrees his
judgments through them, as the Psalmist says, ``thou hast sat on the throne,
who judgest justice'' (Ps 9:5). \emph{Cherubim} means fullness of knowledge,
and those loftier hosts are so called because they are filled with a
knowledge the more perfect the nearer they contemplate the brightness of God.
\emph{Seraphim} means burning or kindling:

\begin{quote}
\latin{Quae, quia ita Deo conjuncta sunt ut inter haec et Deum nulli alii
spiritus intersint, tanto magis ardent, quanto hunc vicinius vident. Quorum
profecto flamma amor est, quia quo subtilius claritatem divinitatis ejus
aspiciunt, eo validius in ejus amore flammescunt.} (\emph{Hom.\ in
Ev.}\ 34.10)
\end{quote}

\noindent They are so joined to God that no other spirits stand between him
and them, and they burn the more, the more nearly they see him. Their flame
is love: the more keenly they gaze on the brightness of his Godhead, the more
strongly they blaze in his love.

Aquinas receives these definitions as the complement of the Dionysian ones.
Dionysius explains the names ``accordingly as they befit the spiritual
perfections they signify,'' while Gregory ``seems to regard more the
exterior ministrations'' (\ST{I q.~108 a.~5}); Aquinas takes from Gregory the
distinction between presiding and dominating (ad~3), the powers who restrain
the hostile spirits (\ST{I q.~109 a.~4}, \emph{sed contra}), and the virtues
through whom miracles are done (\ST{I q.~110 a.~4} obj.~1), and he assigns
the guardianship of each man to the lowest order ``whose place it is,
according to Gregory, to announce the `lesser things'\,'' (\ST{I q.~113 a.~3};
\S\ref{sec:q113}).

\subsubsection{The elect in the choirs of the angels}

The homily turns from the choirs to the hearers. \latin{Sed quid prodest nos
de angelicis spiritibus ista perstringere, si non studeamus haec etiam ad
nostros profectus congrua consideratione derivare?}---what does it profit us
to touch on these things about the angelic spirits unless we take care to
draw them, by fitting reflection, to our own advance? The heavenly city is
made of angels and men, and Gregory holds that as many of the human race
ascend to it as elect angels remained there, reading Moses' song in the
older Latin form, \latin{Statuit terminos gentium secundum numerum angelorum
Dei}, he set the bounds of the nations according to the number of the angels
of God (Deut 32:8), where the Douay, with the Vulgate, reads
``according to the number of the children of Israel.'' Those who return to
the heavenly country must therefore imitate something of the hosts to which
they return:

\begin{quote}
\latin{Distincte namque conversationes hominum singulorum agminum ordinibus
congruunt, et in eorum sortem per conversationis similitudinem
deputantur.} (\emph{Hom.\ in Ev.}\ 34.11)
\end{quote}

\noindent The different ways of life of men correspond to the orders of the
several hosts, and men are assigned to their lot by the likeness of their
lives. Gregory walks up the ladder of the choirs. Those who grasp little but
faithfully announce that little to their brethren run into the number of the
angels; those who grasp and announce the highest heavenly secrets, into the
archangels; those who work signs, into the virtues; those who drive evil
spirits from the possessed by prayer, into the powers; those who surpass the
merits of the elect and rule their brethren, into the principalities; those
who so master every vice that by the right of purity they are called gods
among men, as Moses was appointed ``the god of Pharao'' (Exod 7:1), into the
dominations; those who examine themselves with watchful care and so receive the
gift to judge others, into the thrones, for the Lord presides in them as on
his seat; those full of the love of God and neighbour, into the cherubim, for
cherubim is the fullness of knowledge, and love ``is the fulfilling of the
law'' (Rom 13:10). And at the summit:

\begin{quote}
\latin{Et sunt nonnulli qui, supernae contemplationis facibus accensi, in
solo conditoris sui desiderio anhelant, nihil jam in hoc mundo cupiunt, solo
aeternitatis amore pascuntur, terrena quaeque abjiciunt, cuncta temporalia
mente transcendunt, amant et ardent, atque in ipso suo ardore requiescunt,
amando ardent, loquendo et alios accendunt, et quos verbo tangunt, ardere
protinus in Dei amore faciunt.} (\emph{Hom.\ in Ev.}\ 34.11)
\end{quote}

\noindent Kindled by the torches of heavenly contemplation, they pant for
their Maker alone, are fed by the love of eternity, love and burn and rest in
their burning, and kindle others by their words. What shall I call them but
seraphim, \latin{quorum cor in igne conversum lucet et urit},
whose heart, turned to fire, gives light and burns (34.11)?

Gregory presses the lesson on his hearers. Let each look within and see
whether he finds among these hosts the lot of his calling. \latin{Vae autem
animae quae in se de his bonis quae enumeravimus minime aliquid
recognoscit}: woe to the soul that recognizes in itself nothing of these
gifts, and a worse woe if it knows itself deprived of them and does not
grieve. But whoever finds lesser gifts in himself must not envy greater gifts
in others, \latin{quia et supernae illae distinctiones beatorum spirituum
ita sunt conditae, ut aliae aliis sint praelatae}, because those heavenly
distinctions of the blessed spirits were so created that some are set above
others (34.12). Aquinas determines the same assumption of men into the choirs
through grace: men cannot pass into the angelic grade of nature, but ``by
the gift of grace men can merit glory in such a degree as to be equal to the
angels, in each of the angelic grades'' (\ST{I q.~108 a.~8}).

\subsubsection{Those who stand and those who are sent}

Gregory now cites Dionysius by name:

\begin{quote}
\latin{Fertur vero Dionysius Areopagita, antiquus videlicet et venerabilis
Pater, dicere quod ex minoribus angelorum agminibus foras ad explendum
ministerium vel visibiliter vel invisibiliter mittuntur~\dots\ Nam
superiora illa agmina ab intimis nunquam recedunt, quoniam ea quae
praeeminent usum exterioris ministerii nequaquam habent.} (\emph{Hom.\ in
Ev.}\ 34.12)
\end{quote}

\noindent The lesser hosts are sent out, visibly or invisibly, to fulfil a
ministry; the higher never withdraw from the inmost presence. Isaias seems to
say the contrary: ``one of the seraphims
flew to me, and in his hand was a live coal, which he had taken with the
tongs off the altar. And he touched my mouth'' (Isa 6:6--7). Gregory answers
that the spirits who are sent take the name of those whose office they
bear: the angel who carries a coal from the altar to burn away the sins of
speech is called a seraph because seraph means burning. Daniel's vision
supports the rule, ``thousands of thousands ministered to him, and ten
thousand times a hundred thousand stood before him'' (Dan 7:10):

\begin{quote}
\latin{Aliud namque est ministrare, aliud assistere, quia hi administrant Deo,
qui et ad nos nuntiando exeunt; assistunt vero qui sic contemplatione intima
perfruuntur, ut ad explenda foras opera minime mittantur.} (\emph{Hom.\ in
Ev.}\ 34.12)
\end{quote}

\noindent To minister is one thing, to stand by is another: those minister
to God who also go out to us bearing messages; those stand by who so enjoy
inward contemplation that they are not sent out to fulfil outward works.

Whether the cherubim and seraphim do by themselves what Scripture attributes
to them, or whether it is done by subject hosts who take the names of the
greater because they come from them, Gregory will not affirm, since open
testimonies do not prove it. \latin{Hoc tamen certissime scimus, quia ad
explendum de supernis ministerium alii spiritus alios mittunt}: this we know
most certainly, that some spirits send others to fulfil a ministry from on
high, as Zacharias testifies, ``the angel that spoke in me went forth, and
another angel went out to meet him. And he said to him: Run, speak to this
young man'' (Zach 2:3--4). When an angel says to an angel, \latin{Curre et
loquere}, one plainly sends another: \latin{Minora vero sunt quae mittuntur, majora
quae mittunt}---the lesser are sent, the greater send (34.13). And those who
are sent never leave the presence they came from:

\begin{quote}
\latin{Et mittuntur igitur, et assistunt, quia etsi circumscriptus est
angelicus spiritus, summus tamen spiritus ipse, qui Deus est,
circumscriptus non est. Angeli itaque et missi, et ante ipsum sunt, quia
quolibet missi veniant, intra ipsum currunt.} (\emph{Hom.\ in
Ev.}\ 34.13)
\end{quote}

\noindent They are sent and they stand by, because though the angelic spirit
is bounded, the supreme Spirit, who is God, is not bounded. So the angels are
both sent and before him, for wherever they are sent they run within him.

The names of the orders are proper and common at once. The thrones are God's
seat, and yet the Psalmist says, ``Thou that sittest upon the cherubims''
(Ps 79[80]:2), because the cherubim are joined to the thrones in the
ranks and the Lord is said to sit on them by the nearness of the
neighbouring host. \latin{Sic quippe in illa summa civitate specialia
quaedam singulorum sunt, ut tamen sint communia omnium}: in that highest city
some gifts belong in a special way to each, and yet are common to all, and
what one has in part another order possesses wholly. The seraphim are named
for burning, \latin{et tamen amore conditoris simul omnes ardent}, and yet
all burn together with the love of their Maker; the cherubim for fullness of
knowledge, \latin{et tamen quis ibi aliquid nesciat ubi ipsum omnes simul
fontem scientiae Deum vident?}---and yet who is ignorant of anything there,
where all together see God, the fountain of knowledge? Each order is called
by the proper name of the gift it has received more fully; and even the
gifts that some have and others cannot have, as the dominations and
principalities are named, \latin{cuncta ibi singulorum sunt, quia per
charitatem spiritus ab alio in aliis habentur}: there all things belong to
each, because through charity what one spirit has is had in the others
(34.14). Aquinas gives the same rule when he explains the names of the
orders: ``in the angelic orders all spiritual perfections are common to all
the angels,'' and each order is named from the perfection that belongs to it
as its property (\ST{I q.~108 a.~5}). Then the preacher breaks off: \latin{Suspiremus ergo ad eos de
quibus loquimur, sed redeamus ad nos. Meminisse etenim debemus quia caro
sumus}---let us sigh for those of whom we speak, but return to ourselves,
for we must remember that we are flesh (34.15).

Aquinas receives the Dionysian rule of missions through this part of the
homily. His \emph{sed contra} on whether all
the angels are sent is ``Gregory says (Hom. xxxiv in Evang.), quoting the
statement of Dionysius (Coel. Hier. xiii), that `the higher ranks fulfil no
exterior service'\,'' (\ST{I q.~112 a.~2}); his first objection to the next
article is Gregory's ``the angels are sent, and assist'' (a.~3 obj.~1); and
he reconciles the two sayings by determining that all the angels see the
divine essence and in that sense all assist, while only the highest behold
the secrets of the divine mysteries in the divine essence itself and in
that respect alone are said to assist (a.~3). The rule that the greater send
the lesser answers to Aquinas's principle that ``inferior things are
administered by the superior'' (\ST{I q.~112 a.~2}; \ST{I q.~106 a.~3};
\S\ref{sec:illumination}; \S\ref{sec:q112}).

\subsubsection{The angels in the mysteries of Christ}

At the Christmas Mass in the basilica of Saint Mary, Gregory asks why an
angel appeared to the watching shepherds and the brightness of God shone
round about them: \latin{nisi quod illi prae caeteris videre sublimia
merentur, qui fidelibus gregibus praeesse sollicite sciunt}---because those
above all others merit to see heavenly things who know how to watch with
care over faithful flocks (\emph{Hom.\ in Ev.}\ 8.1). The angels' song
declares a peace made between heaven and earth:

\begin{quote}
\latin{Prius quippe quam Redemptor noster nasceretur per carnem, discordiam
cum angelis habuimus, a quorum claritate atque munditia per primae culpae
meritum, per quotidiana delicta longe distabamus. Quia enim peccando
extranei eramus a Deo, extraneos nos a suo consortio deputabant angeli cives
Dei. Sed quia nos cognovimus Regem nostrum, recognoverunt nos angeli cives
suos.} (\emph{Hom.\ in Ev.}\ 8.2)
\end{quote}

\noindent Before the Redeemer's birth we were at discord with the angels,
far from their brightness by the first fault and by daily sins, and they
counted us strangers; since we have acknowledged our King, they acknowledge
us as fellow citizens. Because the King of heaven took the earth of our
flesh, the angelic height no longer despises our weakness;
\latin{et quos prius infirmos abjectosque despexerant, jam socios
venerantur}---those whom they once looked down upon as weak and abject they
now honour as companions. Gregory finds the sign of this in Scripture. Lot
and Josue adored angels and were not forbidden (Gen 19:1; Jos 5:15), but when
John in the Apocalypse would adore an angel, the angel restrained him:
``See thou do it not. For I am thy fellow servant, and of thy brethren''
(Apoc 22:9). Why do angels receive adoration in silence before the
Redeemer's coming and refuse it afterwards,

\begin{quote}
\latin{nisi quod naturam nostram, quam prius despexerant, postquam hanc super
se assumptam conspiciunt, substratam sibi videre pertimescunt?}
(\emph{Hom.\ in Ev.}\ 8.2)
\end{quote}

\noindent except that, seeing our nature taken up above themselves, they
fear to see it laid beneath them, and do not disdain man as companion when
they adore the God-man above them? Therefore, Gregory concludes, let no
uncleanness defile us, who in God's eternal foreknowledge are citizens of God and equal to his
angels: \latin{Vindicemus moribus dignitatem nostram}---let us uphold our
dignity by our conduct (8.2). The English Catholic Bible carries the same
reading into its note on the angel's refusal in Apoc 19:10, where ``St.\
Gregory (Hom. 8, in Evang.)'' is said to hold that the angel refused ``in
consideration of the dignity to which our human nature had been raised, by
the incarnation of the Son of God'' (Douay--Rheims, Challoner's note on Apoc
19:10).

On the Epiphany Gregory asks why an angel announced the birth to the
shepherds while a star led the Magi: \latin{Quia videlicet Judaeis, tanquam
ratione utentibus, rationale animal, id est angelus, praedicare
debuit}---to the Jews, who used reason, a rational creature, an angel, had
to preach, while the Gentiles, who knew not how to use reason, were led to the
Lord by signs (\emph{Hom.\ in Ev.}\ 10.1). Aquinas meets the phrase
``rational animal'' as an objection for bodies natural to the angels and
answers that ``Gregory calls the angel a rational animal metaphorically, on
account of the likeness to the rational nature'' (\ST{I q.~51 a.~1} ad~2;
\S\ref{sec:q51}).

On the Ascension Gregory places man in the scale of creation: stones exist
but do not live, trees live but do not feel, beasts live and feel but do not
discern, \latin{Angeli etenim sunt, vivunt, sentiunt, et discernunt}---angels
exist, live, perceive, and discern---and man has something of every creature,
\latin{esse cum lapidibus, vivere cum arboribus, sentire cum animalibus,
intelligere cum angelis}: being with stones, life with trees, sense with
animals, understanding with angels (\emph{Hom.\ in Ev.}\ 29.2). Aquinas
makes the last words the \emph{sed contra} of his article on whether the
angels have only intellectual knowledge (\ST{I q.~54 a.~5}). At the Ascension the angels themselves keep the feast:

\begin{quote}
\latin{Quid est ergo quod nato Domino, non in albis vestibus, ascendente
autem Domino, in albis vestibus angeli apparent, nisi quod tunc magna
solemnitas angelis facta est, cum coelum Deus homo penetravit?}
(\emph{Hom.\ in Ev.}\ 29.9)
\end{quote}

\noindent Why do the angels appear in white garments at the Ascension and not
at the Nativity, but that a great feast was made for the angels when God, made
man, entered heaven? At the Nativity the Godhead was seen humbled, at the
Ascension the manhood was exalted, and white garments befit exaltation.

\subsubsection{The \emph{Morals}: the angels who stand and are sent}

The second book of the \emph{Morals} opens on the scene in heaven: ``Now on
a certain day, when the sons of God came to stand before the Lord, Satan also
was present among them'' (Job 1:6). The sons of God are the elect angels, and
Gregory asks how they can come to stand before God when they ``ever behold
and ever stand nigh,'' and yet are ``ministering spirits, sent forth'' (Heb
1:14). The
answer lies in the subtlety of the angelic nature:

\begin{quote}
For they never so go forth apart from the vision of God, as to be deprived
of the joys of interior contemplation; for if when they went forth they lost
the vision of the Creator, they could neither have raised up the fallen, nor
announced the truth to those in ignorance; and that fount of light, which by
departing they were themselves deprived of, they could in no wise proffer to
the blind. (\emph{Moralia} II.3.3)
\end{quote}

\noindent The spirits of angels ``are indeed bounded by space,'' yet they
contemplate the very source of knowledge: ``For of those things which are
capable of being known, what is there that they know not, who know Him, to
Whom all things are known?'' Their knowledge is vast beside ours and little
beside God's, as ``their very spirits in comparison indeed with our bodies
are spirits, but being compared with the Supreme and Incomprehensible Spirit,
they are Body.'' They are sent because they are bounded, and never depart
because they abide in contemplation: ``Thus they at the same time always behold the
Father's face, and yet come to us; because they both go forth to us in a
spiritual presence, and yet keep themselves there, whence they had gone out,
by virtue of interior contemplation'' (II.3.3). Aquinas quotes Gregory's
teaching here in determining that all the angels, even those who minister, assist before
God in the vision of his essence (\ST{I q.~112 a.~3}).

Satan was present among them because, though he lost his blessed estate, ``he
did not part with a nature like to theirs, and though his deserts sink him,
he is lifted up by the properties of his subtle nature'' (II.4.4). He came
before the Lord, but Scripture does not say that he saw the Lord: ``For he
came to be seen, and not to see. He was in the Lord's sight, but the Lord was
not in his sight; as when a blind man stands in the sun, he is himself
bathed indeed in the rays of light, yet he sees nothing of the light, by
which he is brightened'' (II.4.5). Aquinas answers with this passage the
objection that Satan assisted before God: he ``is not described as having
assisted, but as present among the assistants'' (\ST{I q.~112 a.~3} ad~3).
The exposition goes on to the ways in which God speaks to the angels and the
angels to God, which Aquinas receives in his question on the speech of the
angels (\S\ref{sec:q107}).

Gregory finds the whole universe ordered by the governance of higher over
lower: ``man is over the beasts, the Angels over man, and the Archangels are
set over the Angels,'' as Daniel's prince of Persia and Zacharias's two angels
show: ``For in the actual ministration of the
holy spirits, if the superior Powers did not direct the inferior, one Angel
would never have learnt from the lips of another what he should say to a
man.'' The Creator holds all things by himself alone, yet ``for the purpose
of constituting the defined order characterizing a universe of beauty, He
rules one part by the governance of another,'' and the angelic spirits are
rightly called kings and counsellors, for ``they `consult' for the spiritual
commonwealth, while they unite us to the kingdom as fellow-heirs with
themselves'' (\emph{Moralia} IV.29.55).

On the angels of the nations Gregory asks how God ``maketh peace in his high
places'' (Job 25:2) when the angel says to Daniel, ``the prince of the
kingdom of the Persians resisted me one and twenty days: and behold Michael,
one of the chief princes, came to help me'' (Dan 10:13). There are ``fixed
charges of the Angels set to superintend the regulating of the several
particular nations,'' and the spirits are said to come against one another
when the deeds of the peoples under them are set against one another: the
prince of Persia withstands because the captive people still has something
the dominion of the Persians must purge, and Michael helps because the
prophet's prayers claim its release:

\begin{quote}
For the lofty Spirits that are princes to those nations never fight in
behalf of those that act unjustly, but justly judge and try their deeds. And
when either the guilt or guiltlessness of each separate nation is brought
into the debate of the Court Above, the ruling Spirit of that nation is said
to have won in the conflict or not to have won; the one identical victory of
all of whom, however, is the Supreme Will of their Maker above them, which
Will whilst they ever have before their eyes, what they have not the power
they have not the mind to obtain. (\emph{Moralia} XVII.12.17)
\end{quote}

\noindent Aquinas follows Gregory here: ``according to Gregory (Moral. xvii),
the prince of the kingdom of Persia was a good angel appointed to the
guardianship of that kingdom,'' and the angels' resistance is their
consulting of the divine wisdom on the contrary merits of the peoples
(\ST{I q.~113 a.~8}; \S\ref{sec:guardians}).

To the next question of Baldad, ``Is there any numbering of his soldiers?''
(Job 25:3), Gregory answers that the number of the spirits above is beyond
the reckoning of human reason, ``infinite and definite,'' numbered by God and
numberless to man. He repeats the distinction of Daniel: ``those Powers stand
before Him without a doubt, which never go forth for the communicating things
to men. But those `minister to' Him, who come for discharging the offices of
bearers of tidings; yet these same beings also, by the act of contemplation,
are not withdrawn from the interior world,'' and those who minister are more
than those who stand (XVII.13.18). Aquinas uses both halves, the definition
of those who assist and the greater number of those who minister (\ST{I q.~112
a.~3}, \emph{sed contra}; a.~4 obj.~2 and ad~2). The angels are rightly
called God's soldiers, ``because we are not unaware that those war against
the powers of the air, which same conflicts however they carry on not by
labour but by authority,'' and to this host, which sang at the Lord's birth
(Luke 2:13), the number of the elect of men is joined (XVII.13.19).

The morning stars who praised God at the founding of the earth (Job 38:7) are
the angels: ``because the nature of rational spirits is believed to have
been created first in time, the Angels are, not improperly, called `morning
stars.'\,'' They praise ``together'' with men because they adapt the words of
their exultation to our redemption, and ``they rejoice that their own number
is filled up'' when they see us admitted; and they are morning stars also
because, sent to exhort men, ``while they announce the coming morn, they
drive away from the hearts of men the darkness of the present life''
(\emph{Moralia} XXVIII.14.34). Aquinas records two opinions of the Fathers
on whether the angels were created before the corporeal world and holds the
more probable to be that they were created together with it (\ST{I q.~61
a.~3}; \S\ref{sec:q61}).

\subsubsection{The \emph{Morals}: the first of the ways of God}

The angelic nature, Gregory teaches, is in itself changeable, and its
steadfastness is the work of love. On Eliphaz's saying, ``they that serve him
are not steadfast, and in his angels he found wickedness'' (Job 4:18), he
writes:

\begin{quote}
For if the essence of the Angels had been strange to the motion of change,
being created well by its Maker, it would never have fallen in the case of
reprobate spirits from the tower of its blessed estate. But Almighty God in a
marvellous manner framed the nature of the highest spiritual existences good,
yet at the same time capable of change; that both they, that refused to
remain, might meet with ruin, and they, that continued in their own state of
creation, might henceforth be stablished therein more worthily in proportion
as it was owing to their own choice, and become so much the more meritorious
in God's sight, as they had staid the motion of their mutability by the
stablishing of the will. (\emph{Moralia} V.38.68)
\end{quote}

\noindent The angelic nature has overcome its mutability ``in that it is bound
by the chains of love to Him, Who is ever the Same'' (V.38.68). Isidore
takes this sentence into his own chapter on the angels, and Aquinas's
questions on the grace and merit of the angels and on their confirmation
answer the same question in the schools' terms (\S\ref{sec:q62}).

Why was fallen man restored and the fallen angel not? God ``had made two
creations to contemplate Himself, viz. the Angelic and the human, but pride
smote both.'' Man bore the clothing and the infirmity of the flesh; the angel
is spirit alone, and ``it was also meet that the apostate Angel should be
driven down to a farther depth,'' since he fell from his steadfastness
bearing no weakness of the flesh. And there is a second reason: ``the Angel fell by his own
wickedness, but the wickedness of another brought man down''
(\emph{Moralia} IV.3.8). Therefore, Gregory adds, the Redeemer took not the
nature of angels but the seed of Abraham (Heb 2:16), that he ``might at once
let go the lost angel, by not taking his nature, and restore man, by taking
his nature in Himself'' (IV.7.12). Aquinas meets the second reason as an
objection to the obstinacy of the demons and answers that man's sinning at
another's suggestion ``is not the whole cause of man's sin being pardonable''
(\ST{I q.~64 a.~2} obj.~4 and ad~4; \S\ref{sec:q64}).

The greatest angelological passage of the \emph{Morals} comes in the
exposition of Behemoth. God says of him, ``He is the beginning of the ways of
God'' (Job 40:14[19]), and Gregory reads it of the first angel:

\begin{quote}
For what do we understand by the `ways' of God, but His doings?~\dots\ And
Behemoth is said to be the chief of the ways of God, because doubtless when
He was performing all the work of creation, He created him first, whom He
made more eminent than the other Angels. (\emph{Moralia} XXXII.23.47)
\end{quote}

\noindent Ezechiel saw the same eminence in the tree of the paradise of God,
whom ``the cedars~\dots\ were not higher'' (Ezek 31:8--9): the cedars, firs,
and planes are ``those bands of heavenly virtues of lofty height, planted in
the verdure of eternal joy,'' which though created lofty were neither
preferred nor equalled to him, and he was beautiful with thick branches
because ``a comeliness, as great as the subject multitude of Angels which
adorned him, rendered him the more beautiful.'' ``And therefore, when
sinning, he was condemned without pardon, because he had been created great
beyond comparison'' (XXXII.23.47). He is the seal of likeness, as in the
homily, for ``from the seal of a ring such a likeness is impressed in image,
as exists in essence in the seal itself,'' and the nine stones of his covering
are again the nine orders:

\begin{quote}
He mentioned nine kinds of stones, doubtless because there are nine orders
of angels. For when in the very words of Scripture, Angels, Archangels,
Thrones, Dominations, Virtues, Princedoms, Powers, Cherubim, and Seraphim,
are plainly spoken of and mentioned, it is shewn how great are the
distinctions of the citizens of heaven. (\emph{Moralia} XXXII.23.48)
\end{quote}

\noindent The order of this enumeration differs from the ascending list of
the homily; Isidore will copy it. Gregory then finds charity in the
goldsmith's work of Ezek 28:13, ``gold the work of thy beauty: and thy pipes
were prepared in the day that thou wast created.'' Precious stones are
pierced so that gold poured into the holes may bind them into one ornament;
``The holes of this stone were prepared then in the day of its creation,
because, namely, he was created capable of love.'' Had he let himself be
penetrated by ``the gold of charity,'' he would still be ``a stone firmly
fixed in the ornament of a king''; ``he therefore fell, because, even though
perforated with the hand of the artificer, he scorned to be bound with the
bands of love.'' The other stones, pierced like him, ``were bound together by
charity mutually penetrating them, and obtained, on his fall, this, as a
gift, that they should now be never loosened by falling from the ornament of
the King.'' He is called a cherub, ``outspread and covering'' (Ezek 28:14),
because ``Cherub is interpreted, `Plenitude of knowledge,' and he is
therefore called a Cherub, because he is not doubted to have surpassed all in
his knowledge,'' and he walked amid the stones of fire because ``he dwelt
amid the hearts of Angels, which were kindled with the fire of love''
(XXXII.23.48). Aquinas answers the objection from Ezechiel's ``cherub''
with the same etymologies: the first angel who sinned ``is called, not a
Seraph, but a Cherub,'' because knowledge ``is compatible with mortal sin''
but the heat of charity is not (\ST{I q.~63 a.~7} ad~1).

``Let man then consider what he deserves for his pride on earth, if even an
Angel, placed above other Angels, is cast down in heaven'' (XXXII.23.49).
Yet the verse adds, ``who made him, he will apply his sword'': the sword of
Behemoth is ``his malice in doing hurt,'' bent back by him ``by Whom he was
created naturally good,'' so that he who could not touch the swine without
leave (Matt 8:31) cannot of his own accord injure men made after the likeness
of God (XXXII.24.50).

The last books of the \emph{Morals} turn to Leviathan, whom Gregory also reads
of the apostate angel. ``When he shall raise him up, the angels shall fear,
and being affrighted shall purify themselves'' (Job 41:16[17]): Gregory hears
in it the confirmation of the good angels at the fall of the proud one.

\begin{quote}
Nor do we give up the sense of its true meaning, if we believe that when this
Leviathan was falling from the height of blessedness, the Elect Angels also
were greatly terrified at his fall, in order that, as the fall of pride was
casting him out from their number, their very fear might give them strength
to stand more firmly. (\emph{Moralia} XXXIV.7.12)
\end{quote}

\noindent They were purified because, when he went forth with his reprobate
hosts, the elect alone remained. God ``converted the lapse of those who fell
to the benefit of those who remain; and the fault of the proud is punished,
by the same means by which the increased merits of the humble Angels were
discovered and confirmed. For on the fall of these, it was granted as a
special gift to those that they should never in any wise fall''
(XXXIV.7.13). The fallen one ``has lost the happiness of eternal felicity,
yet he has not lost the greatness of his nature,'' by which he still
surpasses all human things, though he is beneath holy men by the baseness of
his deserts (XXXIV.20.39). Leviathan ``was made to fear no one'' (Job
41:24[25]), and Gregory lays bare the root of his sin:

\begin{quote}
He was indeed so made by nature, as to be bound to feel a chaste fear for his
Creator; that is to say, with a subdued and fearless fear, not with the fear
which love casts out, but with the fear which remains for ever and ever, that
is, which love begets.~\dots\ But this Leviathan, beholding the height of His
loftiness, aimed at the privilege of the fatal liberty of ruling over others,
and being subject to no one, saying, I will ascend above the height of the
clouds, and I will be like the Most High.~\dots\ For having left that First
Cause, to Whom he was bound to adhere, he aimed at being, in a sense, his own
first cause. (\emph{Moralia} XXXIV.21.40)
\end{quote}

\noindent ``For him, whom a slavery akin to freedom exalted, a slavish freedom
cast down'' (XXXIV.21.40). And the last word of the Lord's speech, ``he is
king over all the children of pride'' (Job 41:25[26]), names the root of all
the rest: ``This Leviathan, in order to fall in all the points mentioned
above, smote himself with pride alone.~\dots\ For it is written, Pride is the
beginning of all sin. For by this he himself fell, by this he overthrew men
who followed him. He assaulted the health of our immortality with the same
weapon as he destroyed the life of his own blessedness'' (XXXIV.23.47;
Ecclus 10:15). Aquinas determines with the same tradition that the first sin
of the angel was pride and that the devil desired to be as God (\ST{I q.~63
aa.~2--3}), and his
account of the confirmation of the good angels and the obstinacy of the
fallen stands in \S\ref{sec:fall}. The English Catholic Bible, at the lament
over the king of Tyre, notes that what is said there ``is commonly understood
of Lucifer, the king over all the children of pride'' (Douay--Rheims,
Challoner's note on Ezek 28:12).

\subsubsection{The \emph{Dialogues}: the angels at the passing of souls}

The fourth book of the \emph{Dialogues} opens with Adam, who in paradise ``by
purity of heart and heavenly vision, was present with the quires of the
blessed Angels,'' and lost that light by his fall; his children hear of the
heavenly country inhabited by the angels and, being carnal, doubt it, like a
child born in prison who doubts his mother's tales of the sun, ``for he were
a very foolish child, that thought his mother lied, when she spake of light in
other places, because himself, where he was, beheld nothing else but the
darkness of the prison'' (\emph{Dial.}\ IV.1). ``Almighty God
created three kinds of spirits having life'': the angels without bodies, the
souls of men with bodies that do not die with them, and the souls of beasts
that die with their bodies (IV.3). God, who is invisible, is served by
invisible creatures, ``his holy Angels, and the souls of just men''; and as
the soul moves the body, so ``nothing can be disposed of in this visible
world, but by another creature which is invisible'' (IV.5). Aquinas makes
this sentence, with Augustine's, the \emph{sed contra} of his article on
whether the corporeal creature is governed by the angels (\ST{I q.~110 a.~1},
citing it as \emph{Dial.}\ iv, 6; \S\ref{sec:government}).

Gregory then relates what many whose eyes of soul were purified ``by sincere
faith and plentiful prayer'' have seen at the departure of souls (IV.7). At
the death of God's servants, he teaches, ``heavenly musick is heard, to the
end that whiles they give willing ear to that melody, the soul may have no
leisure to feel, when it departeth from the body'': the paralytic Servulus,
singing hymns with the strangers he lodged, cried out, ``Do ye not hear the
great and wonderful musick which is in heaven?'' and died listening, and a
fragrance filled the place until his burial (IV.14). Round the bed of the
paralysed nun Romula there came a light from heaven that filled the cell,
the sound ``as it were of many that came in,'' and on the fourth night
``two quires singing before the door without,'' and ``whiles these heavenly
funerals were in celebrating before the cell door, that holy soul departed
this life, and was carried in that manner up into heaven'' (IV.15). At the
death of the abbot Stephen of Rieti, ``some of them beheld Angels coming in,
but yet were not able to tell it unto others then present,'' and a great
fear fell upon all, ``by which they perceived of what power he was, that
received his soul going out of this world'' (IV.19).

Gregory relates visions of the judgment of souls in which the angels defend
and the demons accuse. The monk Peter, who died and returned to life, was
being carried to be thrown into the fire of hell when ``suddenly an Angel in
a beautiful attire appeared, who would not suffer him to be cast into those
torments,'' and sent him back with the charge, ``Go thy way back again, and
hereafter carefully look unto thyself, how thou leadest thy life''
(IV.36). A soldier who lay as dead in the pestilence saw a certain Stephen
slip on a bridge over a black and stinking river, and he ``was of certain
terrible men, that rose out of the river, drawn by the legs downward: and by
certain other white and beautiful persons, he was by the arms pulled
upward''; Gregory reads it that the sins of the flesh drew him down and his almsgiving
drew him up, and which prevailed is known to the Judge alone (IV.36).

The fallen angels suffer the fire though they have no bodies, and Gregory
argues from them to the souls of the damned before the resurrection: ``If,
then, the devil and his angels, though without bodies, shall be tormented
with corporal fire, what marvel is it that the souls after their departure,
and before they be united again to their bodies, may in like manner suffer
corporal torments?'' (IV.29; Matt 25:41; \S\ref{sec:q64}). And the book closes
with the angels at the altar:

\begin{quote}
for what right believing Christian can doubt, that in the very hour of the
sacrifice, at the words of the Priest, the heavens be opened, and the quires
of Angels are present in that mystery of Jesus Christ; that high things are
accompanied with low, and earthly joined to heavenly, and that one thing is
made of visible and invisible? (\emph{Dial.}\ IV.58)
\end{quote}

\noindent The Church's liturgy joins with the angels to adore the thrice-holy
God and prays in its funeral rite that the angels lead the departed into
paradise (CCC~335; \S\ref{sec:liturgy}).

\sectionguard
\subsection{Isidore of Seville}\label{sec:fa-isidore}

Isidore, younger brother of the Leander to whom Gregory addressed the
\emph{Morals}, succeeded him as bishop of Seville. He treats the angels in two
works that the Latin schools used as handbooks: the seventh book of the
\emph{Etymologies}, on God, the angels, and the saints, and
the first book of the \emph{Sentences}, cited in the Middle Ages as \emph{De
summo bono}, under which title Aquinas quotes it. His chapters are made
largely from Gregory's homily and \emph{Morals}, with Augustine behind them,
set out as definitions and short sentences. Bede in his last days was
translating ``some selections from the books of Bishop Isidore'' for his
pupils, ``saying, `I would not have my boys read a lie, nor labour herein
without profit after my death'\,'' (Cuthbert, \emph{Letter on the Death of
Bede}, in Bede, \emph{Ecclesiastical History}, ed.\ Sellar).

\subsubsection{The \emph{Etymologies}: the names of the angels}

The chapter \emph{De angelis} begins with the word itself:

\begin{quote}
\latin{Angeli Graece vocantur, Hebraice malachoth, Latine vero nuntii
interpretantur, ab eo quod Domini voluntatem populis nuntiant. Angelorum
autem vocabulum officii nomen est, non naturae. Semper enim spiritus sunt,
sed cum mittuntur, vocantur angeli.} (\emph{Etym.}\ VII.5.1--2)
\end{quote}

\noindent The name \latin{angeli} is Greek; in Hebrew they are
\latin{malachoth}, and in Latin \latin{nuntii}, messengers, because they announce the
Lord's will to the peoples; and the word is the name of an office, not of a
nature, for they are always spirits, but are called angels when they are
sent. Painters give them wings by an artist's licence, to signify their
swift passage through all things, as Scripture says of God, ``Who walkest
upon the wings of the winds'' (Ps 103[104]:3; VII.5.3). Then comes the
count: \latin{Novem autem esse ordines angelorum sacrae Scripturae
testantur, id est angeli, archangeli, throni, dominationes, virtutes,
principatus, potestates, cherubim et seraphim} (VII.5.4)---the list, in the
same sequence, of the \emph{Morals} on Behemoth (XXXII.23.48).

The exposition of the names then follows the homily. The archangels are the
chief messengers, and Isidore adds a Greek etymology of his own:
\latin{Archangeli dicti eo quod primatum teneant inter angelos; ARCHOS enim
Graece, Latine princeps interpretatur. Sunt enim duces et principes, sub
quorum ordine unicuique angelorum officia deputata sunt}---they are called
archangels because they hold the first place among the angels; they are
leaders and princes, under whose order the offices are allotted to each
angel (VII.5.6). Zacharias's two angels prove that archangels preside over
angels, and Isidore gives Gregory's reasoning from the \emph{Morals} almost
word for word: \latin{Si enim in ipsis officiis angelorum nequaquam
potestates superiores inferiores disponerent, nullo modo hoc, quod homini
diceret angelus, ab angelo cognovisset}---if the higher powers did not order
the lower in the very offices of the angels, an angel would never have
learned from an angel what he was to say to a man (VII.5.7--8; cf.\
\emph{Moralia} IV.29.55).

Certain archangels are called by proper names so that the names may show
what they can do in their work (VII.5.9). Gabriel, the strength of God, is
sent where divine power or strength is shown, and so he came to Mary
\latin{eo tempore, quo erat Dominus nasciturus et triumphaturus de mundo},
at the time when the Lord was to be born and to triumph over the world, to
announce him \latin{qui ad debellandas aerias potestates humilis venire
dignatus est} (VII.5.10--11). Michael, \latin{Qui sicut Deus}, Who is as God, is sent when a
work of wonderful power is done, \latin{Et ex ipso opere nomen est eius, quia
nemo valet facere quod facere potest Deus}---his name comes from the work
itself, because no one can do what God can do (VII.5.12). Raphael, healing
or medicine of God, is sent wherever a work of healing is needed, and so he
restored Tobias's sight (VII.5.13--14). Isidore adds a fourth, Uriel, fire
of God, after the fire in the bush (VII.5.15); the Church's
discipline discourages names other than Gabriel, Raphael, and Michael, whose
names Holy Scripture contains (\emph{Directory on Popular Piety}~217).

The offices of the higher orders follow Gregory's sequence and definitions,
from the virtues who work signs to the seraphim who burn \latin{quia inter eos
et Deum nulli angeli consistunt}, because no angels stand between them and
God; Isidore sets Daniel's distinction of those who minister and those who
stand by under the principalities, and identifies the cherubim with the two
figures of metal over the propitiatory of the ark, made \latin{propter
significandam angelorum praesentiam, in quorum medio ostenditur Deus}, to
signify the presence of the angels in whose midst God is shown; and the names
are proper and common at once, as Gregory taught (VII.5.16--27).

Isidore then adds three sentences of his own gathering. The offices of the
angels were merited at the beginning: \latin{Vnicuique enim, sicut
praedictum est, propria officia sunt iniuncta, quae promeruisse eos in mundi
constat exordio}---to each, as was said, its proper offices were assigned,
which they are known to have merited at the world's beginning (VII.5.28).
The angels preside over places and men, as Daniel's prince of Persia shows,
\latin{Vnde apparet nullum esse locum cui angeli non praesint. Praesunt enim
et auspiciis operum omnium}---there is no place over which angels do not
preside, and they preside over the beginnings of every work (VII.5.28--29).
And the whole distinction of the orders is that of the angels who stood:

\begin{quote}
\latin{Nimirum ostendentis quod post ruinam angelorum malorum hi, qui
permanserunt, firmitatem aeternae perseverantiae consecuti sunt, nullo iam
lapsu aversi, nulla superbia cadentes, sed firmiter in Dei amore et
contemplatione manentes, nihil aliud dulce habent nisi eum a quo creati
sunt.} (\emph{Etym.}\ VII.5.31)
\end{quote}

\noindent After the apostate angels fell, those who remained were made firm
in the perseverance of eternal blessedness, and so Scripture, after the
creation of heaven, says, \latin{Fiat firmamentum, et vocatum est firmamentum
caelum} (Gen 1:6, 8): it shows that after the ruin of the evil
angels those who remained obtained the firmness of eternal perseverance,
turned away now by no lapse, falling by no pride, but abiding firmly in the
love and contemplation of God, and finding nothing sweet but him by whom they
were made (VII.5.30--31). The chapter closes with Isaias's two seraphim read
as the two Testaments, whose threefold ``Holy'' \latin{Trinitatis in una
divinitate demonstrat mysterium}, shows the mystery of the Trinity in one
Godhead (VII.5.32--33).

\subsubsection{The \emph{Sentences}: the nature, fall, and ministry of the
angels}

The tenth chapter of the first book of the \emph{Sentences} gathers the
teaching into twenty-nine short sentences, drawn chiefly from Gregory's
homily and \emph{Morals}.\footnote{The \emph{Sentences} are cited from
Arevalo's text in Migne, PL~83, 553--560, by chapter and sentence.} It
opens with the name of office (I.10.1) and at once states the ground of the
angels' steadfastness: their nature is changeable, \latin{sed facit eos
incorruptos charitas sempiterna}---but eternal charity makes them
incorruptible. \latin{Gratia dicimus, non natura esse incommutabiles
angelos. Nam si natura incommutabiles essent, diabolus non utique
cecidisset}: we say the angels are unchangeable by grace, not by nature, for
if they were unchangeable by nature the devil would not have fallen. The
contemplation of the Creator supports their changeable nature, and the
apostate angel lost it \latin{dum fortitudinem suam, non a Deo sed a se
voluit custodiri}, when he willed that his strength be kept not by God but
by himself (I.10.2; cf.\ \emph{Moralia} V.38.68).

Isidore places the angels first among creatures. \latin{Ante omnem creaturam
angeli facti sunt, dum dictum est Fiat lux}: the angels were made before
every creature, when it was said \latin{Fiat lux} (Gen 1:3); they are
called light by participating in the eternal light, and wisdom by cleaving to
the unbegotten Wisdom (I.10.3). They were created before the whole creation
of the world, and the devil before all the angels, as it is written, ``He is
the beginning of the ways of God''; he was created first \latin{ordinis
praelatione, non temporis quantitate}, by pre-eminence of rank, not by length
of time (I.10.4). Aquinas, recording both opinions of the Fathers, holds the
more probable to be that the angels were created together with the corporeal
world (\ST{I q.~61 a.~3}). The devil's primacy was the source of the
confidence by which he fell beyond repair; he was the seal of the likeness of
God; as soon as he was made he broke out into pride and, made in the truth,
did not stand in it (I.10.5--7; John 8:44). Man and the devil fell by one
lapse of pride, but man returns by penance, while the apostate angels have no
pardon because no infirmity of the flesh weighed them down when they sinned
(I.10.8--11).

The good angels, after the fall of the apostates, were made firm in eternal
blessedness (I.10.12), and \latin{Bonorum angelorum numerus, qui post ruinam
angelorum malorum est diminutus, ex numero electorum hominum supplebitur, qui
numerus soli Deo est cognitus}---the number of the good angels, lessened by
the ruin of the evil, will be filled up from the number of elect men, a number
known to God alone (I.10.13). Among the angels there is a difference of
powers:

\begin{quote}
\latin{Inter angelos distantia potestatum est, et pro graduum dignitate
ministeria eisdem sunt distributa; aliique aliis praeferuntur, tam culmine
potestatis quam scientia virtutis; subministrant igitur alii aliorum
praeceptis, atque obediunt jussis.} (\emph{Sent.}\ I.10.14)
\end{quote}

\noindent There is among the angels a distance of powers, and ministries are
distributed to them according to the dignity of their grades; some are set
above others both in height of power and in knowledge of virtue; some
therefore serve under the commands of others and obey their orders, as in
Zacharias an angel sends an angel and tells him what to announce. The nine
orders follow in the list of the \emph{Etymologies}, with Ezechiel's nine
stones, which the apostate angel wore before his fall as an ornament fixed on
his garment, and by comparison with which, seeing himself brighter than all,
he swelled at once and lifted his heart to pride (I.10.15). \latin{Angeli
semper in Deo gaudent, non in se}: the angels always rejoice in God, not in
themselves, and the devil is evil because he sought not the things of God but
his own (I.10.16).

On the knowledge of spirits Isidore writes that the angels know all things
in the Word of God before they come to pass, and know by God's revelation
what for men is still to come; and of the fallen: \latin{Praevaricatores
angeli, etiam sanctitate amissa, non tamen amiserunt vivacem creaturae
angelicae sensum. Triplici enim modo praescientiae acumine vigent, id est,
subtilitate naturae, experientia temporum, revelatione superiorum
potestatum}---the fallen angels, though they lost their holiness, did not
lose the keen sense of the angelic creature; they are strong in the edge of
foreknowledge in three ways, by subtlety of nature, by experience of times,
and by revelation from higher powers (I.10.17). Aquinas twice meets the
``experience'' of this sentence as an objection and answers that experience
``is affirmed of angels and demons simply by way of similitude, forasmuch as
they know sensible things which are present, yet without any discursion''
(\ST{I q.~58 a.~3} obj.~3 and ad~3; \ST{I q.~54 a.~5} obj.~2;
\S\ref{sec:knowledge}).

The ministries follow. Whenever God is angry with the world and strikes it
with any scourge, the apostate angels are sent to the ministry of vengeance,
yet restrained by divine power so that they do not harm as much as they
desire; the good angels are deputed to the ministry of human salvation, to
administer the cares of the world and rule all things at God's command, as
the Apostle says, ``Are they not all ministering spirits?'' (I.10.18; Heb
1:14). \latin{Angeli corpora in quibus hominibus apparent, de superno aere
sumunt}: the bodies in which the angels appear to men they take from the
upper air, putting on a solid form from the heavenly element so that they may
be shown more plainly to human eyes (I.10.19), the teaching Aquinas gives
when he says that ``the angels assume bodies of air, condensing it by the
Divine power'' (\ST{I q.~51 a.~2} ad~3; \S\ref{sec:bodies}). \latin{Singulae
gentes praepositos angelos habere creduntur}: each nation is believed to
have angels set over it, as the angel's words to Daniel about the prince of
Persia and Michael show (I.10.20); and all men are shown to have angels, by
the Lord's word, ``their angels in heaven always see the face of my Father,''
and by
the apostles' answer to Rhode, ``It is his angel'' (I.10.21; Matt 18:10;
Acts 12:15).

How can the angels always see the Father's face, and yet desire to look upon
him (1~Pet 1:12)? Isidore answers that they both see and desire
to see, possess and long to possess; if they desired without fulfilment,
their desire would be a penal necessity, and if they were sated, satiety
would bring weariness; therefore \latin{desiderant sine labore, et satiantur
sine fastidio}---they desire without toil and are filled without weariness
(I.10.22--24). Wherever Scripture puts ``angel'' for God, it is the Son who is
understood, by reason of the dispensation of the Incarnation (I.10.25). And
the chapter ends with Gregory's Christmas homily: \latin{Ante dominicae
incarnationis adventum discordia inter angelos et homines fuit. Veniens
autem Christus pacem in se, et angelis, et hominibus fecit}---before the
coming of the Lord's Incarnation there was discord between angels and men,
but Christ in his coming made peace in himself for angels and for men
(I.10.26). The sign of that discord is that in the Old Testament men greeted
angels and were not greeted in return, while in the New the angel not only
receives John with reverence but forbids him to adore (I.10.27--28); now man
is received by God and reverently greeted by the angel, as Gabriel greeted
Mary and as the angel told John, ``I am thy fellow servant'' (I.10.29; Apoc
22:9).

\sectionguard
\subsection{The Venerable Bede}\label{sec:fa-bede}

Bede, monk of Wearmouth and Jarrow, finished the \emph{Ecclesiastical History of the
English People} in 731 and died on the eve of the Ascension in 735. In his
commentaries he hands on Gregory's teaching in Gregory's own words, and in
his \emph{History} he records the visions in which the English Church saw the
angels meet the dying, carry the souls of the saints, and contend with the
demons for the souls of sinners. He owed the faith of his people, as he tells
it, to Gregory's love for the English, and he tells the tradition of its
beginning. Gregory saw fair-haired boys for sale in the market at Rome and was
told they were Angles: `` `Right,' said he, `for they have an angelic face,
and it is meet that such should be co-heirs with the Angels in heaven'\,''
(\emph{Hist.\ eccl.}\ II.1). Bede relates it as ``the story of the blessed
Gregory, handed down to us by the tradition of our ancestors.''

\subsubsection{The commentaries: Gabriel and the heavenly host}

On the Annunciation Bede's commentary on Luke reproduces the homily on the
drachma: \latin{Idcirco angeli privatis nominibus censentur, ut signetur per
vocabula etiam in operatione quid valeant}, the angels bear proper names so
that their names may show what they can do in their work; they do not
receive names in the holy city, but draw them from their ministries when
they come to us; and Gabriel, the strength of God, is sent to Mary
\latin{qui ad debellandas aereas potestates humilis apparere dignatus est}
(\emph{In Luc.}\ I, on Luke 1:26; PL~92, 315--316; cf.\ \emph{Hom.\ in
Ev.}\ 34.8--9). His homily for the feast of the Annunciation sets the
angel over against the serpent:

\begin{quote}
\latin{Aptum profecto humanae restaurationis principium, ut angelus a Deo
mitteretur ad virginem partu consecrandam divino, quia prima perditionis
humanae fuit causa, cum serpens a diabolo mittebatur ad mulierem spiritu
superbiae decipiendam~\dots\ Illa a diabolo seducta per serpentem, viro
gustum necis obtulit; haec a Deo edocta per angelum, mundo auctorem salutis
edidit.} (\emph{Hom.}\ I.1; PL~94, 9)
\end{quote}

\noindent As the serpent was sent by the devil to deceive a woman by the
spirit of pride, so the angel was sent by God to the Virgin; the one woman,
seduced through the serpent, offered her husband the taste of death, the
other, taught through the angel, brought forth the Author of salvation. Bede
notes that angels who appear to men are rarely named, and that when they
are, the name shows what they come to minister; Gabriel is the strength of
God and rightly shines with that name, \latin{quia nascituro in carne Deo
testimonium perhibet}, because he bears witness to God about to be born in
the flesh. And Mary \latin{jure angelico aspectu simul et affatu meruit
perfrui, quae angelicam studebat vitam imitari}---rightly merited to enjoy
both the sight and the speech of an angel, since she strove to imitate the
angelic life (PL~94, 9--10).

On the multitude of the heavenly army at Bethlehem (Luke 2:13) Bede joins
the homily, the \emph{Morals}, and the Christmas sermon into one exposition:

\begin{quote}
\latin{Et bene chorus adveniens angelorum militiae coelestis vocabulum
accepit, qui et duci illo potenti in praelio, qui ad debellandas aereas
potestates apparuit, humiliter obsecundat, et ipse potestates easdem
contrarias, ne mortales tantum tentare valeant quantum volunt, fortiter
armis coelestibus proturbat.} (\emph{In Luc.}\ I, on Luke 2:13--14; PL~92,
333)
\end{quote}

\noindent The choir is rightly called an army, for it serves the Captain
mighty in battle and itself drives back the hostile powers with heavenly
arms, lest they tempt mortals as much as they will. As a good emperor
fortifies every place against the enemy's coming, \latin{ita et Deus, quoniam immundi spiritus ad pacis eversionem
ubique versantur, ad tutelam nostram constituit exercitus angelorum, quorum
praesentia et daemonum confringatur audacia, et nobis pacis gratia
ministretur}: so God, since unclean spirits roam everywhere to overthrow
peace, has set armies of angels to guard us, by whose presence the boldness
of the demons is broken and the grace of peace is ministered to us. The
angels glorify God incarnate for our redemption \latin{quia dum nos
conspiciunt recipi, suum gaudent numerum repleri}, because when they see us
received they rejoice that their number is filled up, and they wish peace to
men \latin{quia quos infirmos prius abjectosque despexerant, nascente in
carne Domino jam socios venerantur}---because those they once looked down
upon as weak and abject they now, the Lord being born in the flesh, honour as
companions (PL~92, 333; cf.\ \emph{Hom.\ in Ev.}\ 8.2 and \emph{Moralia}
XXVIII.14.34).

\subsubsection{The \emph{Explanation of the Apocalypse}: Michael and the
Church}

In his commentary on the Apocalypse Bede expounds the war in heaven (Apoc
12:7):

\begin{quote}
The heaven signifies the Church, in which he says that Michael, with his
angels, fights against the devil, for that, according to the will of God, he
contends for the Church in her sojourning, by praying and ministering help;
of whom Daniel also said, that he would come to the aid of the Church in the
last and most grievous affliction; from which they suppose that Antichrist
is to be slain by him. And they are said to be his angels in the same way
that our angels also are. (\emph{Expl.\ Apoc.}\ II, on 12:7)
\end{quote}

\noindent Daniel's word is that Michael ``shall rise up, the great prince,
who standeth for the children of thy people'' (Dan 12:1), and the angels are
called Michael's as ours are called ours, the angels of those whose citizens
they are (Matt 18:10). The angels of Satan are not only those
like him in nature and will, ``but men who are entangled in their snares are
also to be understood.'' When the loud voice says ``the accuser of our
brethren is cast forth'' (Apoc 12:10), ``The angels express joy at the
salvation of their brethren, that is, of those who will become citizens, but
who now are strangers''; and the call to rejoice, ``O heavens, and you that
dwell therein'' (Apoc 12:12), is addressed to angels and holy men together,
``since both men are joined with angels, and angels minister to man's nature
in Christ'' (\emph{Expl.\ Apoc.}\ II, on 12:10--12).

Bede keeps Gregory's reading of the angel who refused John's adoration: ``After
that the Lord Jesus Christ raised the person of man, which He assumed, above
the heavens, the angel feared to be worshipped by man, namely, as worshipping
the God-man above himself. Yet we read of this having been done before the
Incarnation of the Lord by men, and not in any wise forbidden by the angels''
(\emph{Expl.\ Apoc.}\ III, on 19:10; cf.\ \emph{Hom.\ in Ev.}\ 8.2).
Aquinas cites Bede elsewhere in the treatise on the angels, for the saying
that ``the devil does not send wicked thoughts, but kindles them'' (\ST{I
q.~111 a.~2} obj.~2).

\subsubsection{The \emph{Ecclesiastical History}: the angels and the dying}

The \emph{History} relates the angels' ministry at the death of the saints of
the English Church. When the virgin Earcongota died in the Frankish
monastery of Brige, the brethren in the other houses
``declared they had then plainly heard choirs of singing angels, and, as it
were, the sound of a multitude entering the monastery,'' and saw ``a great
light coming down from heaven, which bore that holy soul, set loose from the
bonds of the flesh, to the eternal joys of the celestial country''
(\emph{Hist.\ eccl.}\ III.8). The monk Owini, working outside while
Bishop Chad prayed alone in his retired dwelling near the church, heard ``a sweet sound
of singing and rejoicing descend from heaven to earth,'' enter the oratory,
and after half an hour return to heaven. Chad summoned the brethren,
announced his death, and to Owini alone explained the song:

\begin{quote}
``If you heard the singing, and know of the coming of the heavenly company, I
command you, in the Name of the Lord, that you tell it not to any before my
death. But in truth they were angelic spirits, who came to call me to my
heavenly reward, which I have always loved and longed after, and they
promised that they would return seven days hence, and take me away with
them.'' (\emph{Hist.\ eccl.}\ IV.3)
\end{quote}

\noindent On the seventh day, having received the Body and Blood of the Lord,
he departed, ``led, as may justly be believed, by the attendant angels,'' and
Egbert in Ireland spoke of one who saw ``the soul of his brother Cedd, with a
company of angels, descending from heaven, who, having taken Ceadda's soul
along with them, returned again to the heavenly kingdom'' (IV.3). When the
abbess Hild died at Whitby, the nun Begu at Hackness, thirteen miles away,
heard the bell that called the sisters to prayer at a death, saw the roof
open and a light pour down, and saw ``the soul of the aforesaid handmaid of
God in that same light, being carried to heaven attended and guided by
angels''; another nun at Whitby itself ``saw her soul ascend to heaven in the
company of angels'' (IV.23).

The visions of the next world show the angels as guides and defenders. The
Irish monk Fursey, in a trance, ``was accounted worthy to behold the sight of
the choirs of angels, and to hear their glad songs of praise,'' among them
the refrain ``The saints shall go from strength to strength''; taken up again
three days later, he saw ``fierce conflicts of evil spirits, who by frequent
accusations wickedly endeavoured to obstruct his journey to heaven; but the
angels protected him, and all their endeavours were in vain.'' His three
angel guides showed him four fires that would consume the world, of
falsehood, covetousness, discord, and ruthlessness. When the fire drew near
and he feared, the angel answered, ``That which you did not kindle will not
burn you''; one angel divided the flames while the other two defended him on
either side; and when the demons accused him over the garment he had taken
from a dying sinner, ``the angel withstood him, saying, `He did not receive
them through avarice, but in order to save his soul'\,'' (\emph{Hist.\
eccl.}\ III.19). A Northumbrian householder, Drythelm, who died and rose
again, was led through the places of the dead by a guide who ``had a
countenance full of light, and shining raiment.'' Left alone at the mouth of
the pit, beset by dark spirits with fiery tongs, he saw ``the brightness of a
star shining amidst the darkness,'' and at the approach of his guide ``all
those evil spirits, that sought to carry me away with their tongs, dispersed
and fled.'' The guide showed him the valley of purgation, where souls who
repented only at death are tried until the judgment, of whom ``many are
succoured before the day of judgement, by the prayers of the living and their
alms and fasting, and more especially by the celebration of Masses''
(V.12).

A Mercian story sets the angels and the demons as witnesses at the
judgment. A king's thegn who put off confession saw two fair youths bring a
very small and beautiful book of his good deeds, and then an army of evil
spirits bring a book ``terrible to behold, of a monstrous size'' with all his
sins; the youths in white said to the demons, ``take him and lead him away,''
and he died in despair. Bede draws the lesson: ``our deeds and thoughts are
not scattered to the winds, but are all kept to be examined by the Supreme
Judge, and will in the end be shown us either by friendly angels or by the
enemy'' (V.13).

The archangel Michael himself appears in the \emph{History}. When Bishop
Wilfrid lay as though dead at Meaux on his way home from Rome, he woke and
told the priest Acca: ``There stood by
me a certain one, glorious in white raiment, and he told me that he was
Michael, the Archangel, and said, `I am sent to call you back from death: for
the Lord has granted you life, through the prayers and tears of your
disciples and brethren, and the intercession of His Blessed Mother Mary, of
perpetual virginity'\,'' (V.19). The archangel who, in Gregory's homily, is
sent whenever a work of wonderful power is done, is sent here at the prayer of
the Church and the intercession of the Mother of God.

\sectionguard
\subsection{Rabanus Maurus and the Carolingian Schools}\label{sec:fa-rabanus}

Rabanus Maurus, abbot of Fulda and archbishop of Mainz, gave the Carolingian
schools an encyclopedia in twenty-two books, \emph{De universo}, and its
chapter \emph{De angelis} shows how the Latin teaching had settled. Rabanus
copies Isidore's chapter from the \emph{Etymologies} in order, from the
Hebrew \latin{malachoth} to the Trinity in the seraphim's cry, and then,
where Isidore ends, he goes on with the words of Gregory's homily:
\latin{Sed quid prodest nos de angelicis spiritibus ista perstringere, si non
studeamus haec etiam ad nostros profectus congrua consideratione
derivare?} The encyclopedia gives the whole ladder of the homily, the
heavenly city built of angels and men and the elect assigned to the choirs
according to their lives, from those who announce little things piously to
their brethren up to those whose hearts, turned to fire, give light and
burn, with Gregory's warning that he who finds lesser gifts in himself must
not envy greater gifts in others (cf.\ \emph{Hom.\ in Ev.}\ 34.11--12).
It closes with the senses in which Scripture uses the word ``angel'': Christ,
the messenger of the Father's will and the ``angel of the testament'' (Mal
3:1); John the Baptist, ``my angel'' sent before the Lord's face; the holy
preachers, as Gregory read Leviathan's verse, ``the angels shall fear''; and
even the devil and the unclean spirits (Gal 1:8; 1~Cor 6:3), \latin{Quod
nomen non per naturam, sed per officium ministrationis sortiti sunt}---a
name they received not by nature but by the office of ministry (\emph{De
universo} I.5; PL~111, 27--32). In the same century the \emph{Celestial
Hierarchy} itself reached the Latin West, first in the version made for
Hilduin of Saint-Denis after the writings came to Louis the Pious in 827, and
then in the translation that John Scotus Eriugena made about 858 at the
instance of Charles the Bald, through which the medieval schools read
Dionysius (\S\ref{sec:ch}). The schools thus received the Dionysian order
beside Gregory's, and Aquinas's comparison of the two is the fruit of that
double inheritance (\S\ref{sec:q108}).

\sectionguard
\subsection{The Latin Fathers after Dionysius and Their Reception}\label{sec:fa-reception}

The teaching of Gregory, Isidore, Bede, and Rabanus, with the place where
Aquinas or the Church receives it, the visions and narratives of the
passing of souls being graded Pious tradition:

\begin{fourstable}{Witness}{Locus}{Teaching}{Received at}
Gregory & \emph{Hom.\ in Ev.}\ 34.3, 6 & The hundred sheep and ten drachmas are angels and men; man was made the tenth so that the number of the elect be complete & \ST{I q.~108 a.~8}\\
Gregory & \emph{Hom.\ in Ev.}\ 34.7 & Nine orders are gathered from Scripture (Eph 1:21; Col 1:16; the prophets); the nine stones of Ezek 28:13 are the nine orders & \ST{I q.~108 a.~6}; Appendix~\ref{app:orderings}\\
Gregory & \emph{Hom.\ in Ev.}\ 34.7; \emph{Moralia} XXXII.23.47--48 & The first angel was the seal of likeness, set over all the hosts and brighter than all; created first, he fell without pardon & \ST{I q.~63 a.~7} s.c.\ and ad~1\\
Gregory & \emph{Hom.\ in Ev.}\ 34.8 & ``Angel'' is the name of an office, not of a nature; angels announce lesser things, archangels the highest & CCC~329; \ST{I q.~108 a.~5}; \ST{I q.~113 a.~3}\\
Gregory & \emph{Hom.\ in Ev.}\ 34.8--9 & Proper names are taken from ministries: Michael, Gabriel, Raphael & \emph{Directory on Popular Piety}~217\\
Gregory & \emph{Hom.\ in Ev.}\ 34.10 & The office of each order: virtues work miracles, powers curb the demons, principalities preside, dominations possess, thrones judge, cherubim know, seraphim burn & \ST{I q.~108 aa.~5--6}; \ST{I q.~109 a.~4} s.c.; \ST{I q.~110 a.~4} obj.~1\\
Gregory & \emph{Hom.\ in Ev.}\ 34.11--12 & The elect are assigned to the choirs by the likeness of their lives & \ST{I q.~108 a.~8}\\
Gregory & \emph{Hom.\ in Ev.}\ 34.12--13 & With Dionysius, the higher hosts are not sent; to minister and to stand by differ; the greater send the lesser; the sent run within God & \ST{I q.~112 aa.~2--3}; \ST{I q.~106 a.~3}\\
Gregory & \emph{Hom.\ in Ev.}\ 34.14 & Each order's gift is proper and common; through charity all things belong to each & \ST{I q.~108 a.~5}\\
Gregory & \emph{Hom.\ in Ev.}\ 8.2 & Through the Incarnation the angels acknowledge men as fellow citizens and refuse their adoration & Challoner's note on Apoc 19:10; CCC~333\\
Gregory & \emph{Hom.\ in Ev.}\ 10.1; 29.2 & The angel is a rational creature; man understands with the angels & \ST{I q.~51 a.~1} ad~2; \ST{I q.~54 a.~5} s.c.\\
Gregory & \emph{Moralia} II.3.3; II.4.4--5 & The sent angels never leave the vision of God; Satan came among them to be seen, not to see & \ST{I q.~112 a.~3} co.\ and ad~3\\
Gregory & \emph{Moralia} IV.29.55 & Man is over the beasts, angels over man, archangels over angels; one part of the universe is ruled by another & \ST{I q.~110 a.~1}; \ST{I q.~112 a.~2}; \ST{I q.~106 a.~3}\\
Gregory & \emph{Moralia} XVII.12.17 & Angels preside over nations and judge their merits; the prince of Persia is a good angel & \ST{I q.~113 a.~8}\\
Gregory & \emph{Moralia} XVII.13.18--19 & The host is numberless to man; those who minister outnumber those who stand; the angels war on the demons by authority & \ST{I q.~112 a.~3} s.c.; a.~4 ad~2; \ST{I q.~50 a.~3}\\
Gregory & \emph{Moralia} XXVIII.14.34 & The angels, created first, are the morning stars; they rejoice that their number is filled up by men & \ST{I q.~61 a.~3} (creation with the corporeal world held more probable)\\
Gregory & \emph{Moralia} V.38.68 & The angelic nature is changeable and made firm by the bond of love & \ST{I q.~62 a.~8}\\
Gregory & \emph{Moralia} IV.3.8 & The angel fell by his own wickedness, man by another's; the angel had no infirmity of flesh & \ST{I q.~64 a.~2} obj.~4 and ad~4\\
Gregory & \emph{Moralia} XXXIV.7.12--13 & The elect angels were terrified at the fall of the proud and confirmed never to fall & \ST{I q.~62 a.~8}\\
Gregory & \emph{Moralia} XXXIV.21.40; 23.47 & The devil refused chaste fear, sought to be his own first cause, and fell by pride, the root of sin & \ST{I q.~63 aa.~2--3}; Challoner's note on Ezek 28:12\\
Gregory & \emph{Dial.}\ IV.5 & Nothing is disposed in the visible world but by the invisible creature & \ST{I q.~110 a.~1} s.c.\\
Gregory & \emph{Dial.}\ IV.14--15, 19, 36 & Angels and heavenly choirs attend the passing of the just; angels and demons contend for souls (Pious tradition) & CCC~335\\
Gregory & \emph{Dial.}\ IV.29 & The devil and his angels, without bodies, are tormented by corporal fire & \ST{I q.~64 a.~4} ad~3\\
Gregory & \emph{Dial.}\ IV.58 & The choirs of angels are present at the Sacrifice & CCC~335\\
Isidore & \emph{Etym.}\ VII.5.1--4, 16--27 & The name of office; nine orders; Gregory's definitions of each & \ST{I q.~108 a.~5}\\
Isidore & \emph{Etym.}\ VII.5.9--15 & Proper names from ministries; Uriel added & \emph{Directory on Popular Piety}~217\\
Isidore & \emph{Etym.}\ VII.5.28--31 & Angels preside over every place; those who stood were made firm after the fall & \ST{I q.~110 a.~1}; \ST{I q.~62 a.~8}\\
Isidore & \emph{Sent.}\ I.10.2--4 & Unchangeable by grace, not nature; created before all, the devil first in rank & \ST{I q.~62 a.~8}; \ST{I q.~61 a.~3} (creation with the corporeal world held more probable)\\
Isidore & \emph{Sent.}\ I.10.13--14 & The lessened number filled up from elect men; grades of power and ministry & \ST{I q.~108 a.~8}; \ST{I q.~106 a.~3}\\
Isidore & \emph{Sent.}\ I.10.17 & The fallen know by subtlety of nature, experience of times, and revelation & \ST{I q.~58 a.~3} obj.~3 and ad~3; \ST{I q.~54 a.~5} obj.~2\\
Isidore & \emph{Sent.}\ I.10.19--21 & Angels take bodies of air; nations and all men have angels & \ST{I q.~51 a.~2} ad~3; \ST{I q.~113 aa.~2, 4}\\
Bede & \emph{In Luc.}\ I (PL~92, 315--316, 333) & Gregory's names of office; the heavenly army guards men and breaks the demons' boldness & \ST{I q.~109 a.~4}; \ST{I q.~113 a.~1}\\
Bede & \emph{Hom.}\ I.1 (PL~94, 9--10) & The angel sent to the Virgin answers the serpent sent to Eve & CCC~332\\
Bede & \emph{Expl.\ Apoc.}\ II, on 12:7--12 & Michael fights for the pilgrim Church by prayer and help; angels rejoice at men's salvation & CCC~335\\
Bede & \emph{Hist.\ eccl.}\ III.8, 19; IV.3, 23; V.12--13, 19 & Angels carry the souls of saints, defend souls from the demons, and bear witness at judgment (Pious tradition) & CCC~336\\
Rabanus & \emph{De universo} I.5 & Isidore's chapter joined to Gregory's ladder of the elect; ``angel'' a name of office in all its scriptural senses & ---\\
\end{fourstable}
